%%% rubric.tex --- Example of using CurVe.

%% Copyright (C) 2000, 2001, 2002, 2003, 2004, 2005, 2010 Didier Verna.

%% Author:        Didier Verna <didier@lrde.epita.fr>
%% Maintainer:    Didier Verna <didier@lrde.epita.fr>
%% Created:       Thu Dec 10 16:04:01 2000
%% Last Revision: Mon Dec  6 11:04:22 2010

%% This file is part of CurVe.

%% CurVe may be distributed and/or modified under the
%% conditions of the LaTeX Project Public License, either version 1.1
%% of this license or (at your option) any later version.
%% The latest version of this license is in
%% http://www.latex-project.org/lppl.txt
%% and version 1.1 or later is part of all distributions of LaTeX
%% version 1999/06/01 or later.

%% CurVe consists of the files listed in the file `README'.


%%% Commentary:

%% Contents management by FCM version 0.1.


%%% Code:
 \begin{rubric}{Cursus}
                                \entry*[2020-2021]
                            Troisième année de doctorat
                    
                                \entry*
                            Membre scientifique de la Casa de
                        Velázquez
                    
                                \entry*[2019-2020]
                            Seconde année de doctorat.
                    
                                \entry*
                            Assistant de recherche (\textit{Visiting Assistant in
                        Research}), Yale University, \textit{Graduate Department
                        of Spanish and Portuguese}, auprès de Jesús R.
                        \textsc{Velasco}, septembre -- début décembre.
                    
                    
                                \entry*[2018-2019]
                            Première année de doctorat sous la direction de
                        Carlos \textsc{Heusch}\href{https://www.idref.fr/055838413}{\textsuperscript{\includegraphics[scale=0.025]{img/idref.png}}}.
                        Laboratoire de rattachement: \href{http://ciham.ish-lyon.cnrs.fr/}{CIHAM} (UMR
                        6848). Co-direction avec Jesús \textsc{
                        Velasco}\href{https://orcid.org/0000-0002-3848-9230}{\includegraphics[scale=0.025]{img/orcid.png}}.
                    
                                \entry*
                            Suivi de l'atelier DTS (\textit{Distributed Text
                        Services}) \href{https://distributed-text-services.github.io/workshops/events/2019-hamburg/}{\enquote{Texts and API}} à Hambourg, du 15 au 16
                        juillet.
                    
                                \entry*
                            Suivi de l'école d'été de l'IRHT \href{https://www.irht.cnrs.fr/?q=fr/agenda/le-livre-medieval-au-regard-des-methodes-quantitatives-0}{\enquote{Le livre médiéval au regard des méthodes
                        quantitatives}} du 25 au 28 juin.
                    
                                \entry*[2017-2018]
                            Préparation de la thèse (4ème année à l'ÉNS de
                        Lyon). 
                    
                                \entry*
                            Suivi d'un atelier en édition numérique à
                        Cambridge (Royaume-Uni) en avril 2018, dans la
                        continuité du DEMM (\textit{Digital Encoding of Medieval
                        Manuscripts}).
                    
                                \entry*
                            Participation à la session 2018 de la formation
                        EDEEN (École d'été en humanités numériques) à Grenoble,
                        organisée par l'Université Grenoble-Alpes et la MSH
                        Alpes (Grenoble, 28 mai-2 juin).
                    
                                \entry*
                            Participation à l'atelier 2018 du consortium
                        CAHIER (Montpellier, 26-29 juin 2018).
                    
                                \entry*[2016-2017]
                            Préparation et obtention du concours de
                        l'Agrégation d'Espagnol. 
                    
                                \entry*[2015-2016]
                             Master II d'études hispanophones du département
                        d'études lusophones et hispanophones de l'ÉNS de Lyon.
                        Suivi des cours du Master \textit{Patrimonio Textual y
                        humanidades digitales} de l’Université de
                        Salamanque dans le cadre du programme Erasmus. Mémoire
                        préparé sous la direction de Carlos \textsc{Heusch}\href{https://www.idref.fr/055838413}{\textsuperscript{\includegraphics[scale=0.025]{img/idref.png}}} et de Francisco
                        \textsc{Bautista}\href{https://orcid.org/0000-0002-2676-0388}{\includegraphics[scale=0.025]{img/orcid.png}}: \textit{La version B du
                        ``Regimiento de los príncipes
                        glosado''}. Note obtenue: 19. 
                    
                                \entry*
                            Suivi de la formation \textit{Digital Encoding of
                        Medieval Manuscripts} (DEMM) financée par l’Union
                        Européenne. Trois semaines de stage à Vienne (Abbaye de
                        Klosterneuburg), Lyon (ÉNS de Lyon), Londres (Queen Mary
                        University). 
                    
                                \entry*[2014-2015]
                            Master I d'études hispanophones au sein du
                        département d'études lusophones et hispanophones de
                        l'ÉNS de Lyon. Spécialisation en philologie castillane
                        et en étude de textes du bas-Moyen Âge. Mémoire préparé
                        sous la direction de Carlos Heusch: \textit{Edición
                        crítica de dos textos de Pedro de Chinchilla}.
                        Note obtenue: 18.
                    
                                \entry*
                            Stage d’ecdotique des Sources chrétiennes (Lyon,
                        février 2015).
                    
                                \entry*[2013-2014]
                            Licence 3, Université Lumière Lyon 2-ÉNS de Lyon. 
                    
                                \entry*
                            Séminaire d’édition de textes de Gisèle
                        \textsc{Besson}, ÉNS Lyon.
                    
                                \entry*
                            Stage d'introduction au Grec ancien (ILOAM,
                        Université de Lausanne, été 2014).
                    
                                \entry*[2011-2013]
                            Classe préparatoire Littéraire, Lycée Joffre
                        (Montpellier). Préparation aux concours des ÉNS (khâgne
                        classique). Admis à l'ÉNS de Lyon.
                    \end{rubric}




%%% rubric.tex ends here

%%% Local Variables:
%%% mode: latex
%%% TeX-master: "raw"
%%% End: